\chapter{روال‌ها}%
\label{ch:functions}

تا اینجا چند بار با \code{روال} روبه‌رو شده‌ایم دستِ‌کم روالِ \code{ریشه}
را در هر برنامه دیده‌ایم. در این فصل روال‌ها را به‌طورِ کامل می‌شناسیم: چرا به
آن‌ها نیاز داریم، چگونه می‌سازیمشان، چگونه به آن‌ها ورودی می‌دهیم و چگونه از
آن‌ها خروجی می‌گیریم.

\section{روال چیست و چرا به آن نیاز داریم؟}

روال، بسته‌ای از دستورهاست که نامی دارد و کاری مشخص انجام می‌دهد. تصور کنید
دستورِ پختِ یک کیک را یک‌بار نوشته‌اید؛ از آن پس هر وقت کیک خواستید، کافی است
بگویید «کیک بپز» لازم نیست هر بار همهٔ مراحل را تکرار کنید. روال همین نقش
را دارد: یک‌بار می‌نویسیدش و بارها \emph{صدایش می‌زنید}.

روال‌ها سه فایدهٔ بزرگ دارند:
\begin{itemize}
  \item \textbf{پرهیز از تکرار:} کاری که چند بار لازم می‌شود را یک‌بار
        می‌نویسید.
  \item \textbf{خوانایی:} برنامه به بخش‌های کوچک و معنادار تقسیم می‌شود.
  \item \textbf{آزمون‌پذیری:} هر روال را جداگانه می‌توان آزمود و درست کرد.
\end{itemize}

\section{ساختنِ نخستین روال}

به این مثالِ واقعی نگاه کنید که دو روالِ کوچک می‌سازد:

\begin{salamfa}
روال توان(پایه : صحیح, نما : صحیح) : صحیح:
    ناپایا نتیجه : صحیح = 1
    تکرار نما:
        نتیجه = نتیجه * پایه
    پایان
    برگشت نتیجه
پایان

روال بیشینه(الف : صحیح, ب : صحیح) : صحیح:
    اگر الف > ب:
        برگشت الف
    پایان
    برگشت ب
پایان

روال ریشه:
    سرچاپ "دو به توان ده:", توان(2, 10)
    سرچاپ "بیشینه:", بیشینه(توان(2, 4), توان(3, 2))
پایان
\end{salamfa}

ساختارِ کلیِ یک روال چنین است:

\begin{center}
\code{روال} \;\textit{نام}\;\code{(}\textit{ورودی‌ها}\code{)}\; \textit{گونهِ خروجی}\; \code{:} \quad\ldots\;بدنه\;\ldots\quad \code{پایان}
\end{center}

بیایید روالِ \code{توان} را بشکافیم:
\begin{itemize}
  \item \code{توان} نامِ روال است.
  \item \code{(پایه صحیح, نما صحیح)} دو \textbf{پارامتر} (ورودی) است: هر
        کدام نامی و گونهی دارند و با کاما جدا شده‌اند.
  \item \code{صحیح} پس از پرانتز، \textbf{گونهِ خروجیِ} روال است؛ یعنی این
        روال یک عددِ صحیح برمی‌گرداند.
  \item درونِ بدنه، توان را با یک حلقه حساب می‌کنیم و با \code{برگشت نتیجه}
        آن را بیرون می‌دهیم.
\end{itemize}

\section{پارامترها: ورودیِ روال}

پارامترها مانندِ متغیرهایی هستند که هنگامِ \emph{صدا زدنِ} روال، مقدار
می‌گیرند. وقتی می‌نویسیم \code{توان(2, 10)}، عددِ \code{2} در \code{پایه} و
عددِ \code{10} در \code{نما} می‌نشیند. ترتیب اهمیت دارد: نخستین مقدار به نخستین
پارامتر می‌رود و الی آخر.

\begin{note}
هر روال می‌تواند هیچ، یک، یا چند پارامتر داشته باشد. روالِ \code{ریشه} معمولاً
بی‌پارامتر است؛ به همین دلیل پس از نامش پرانتزِ خالی نمی‌گذاریم و مستقیم
\code{روال ریشه:} می‌نویسیم.
\end{note}

\section{مقدارهای پیش‌فرضِ پارامتر}

یک پارامتر می‌تواند با \code{=} یک \textbf{مقدارِ پیش‌فرض} داشته باشد. هرگاه
فراخواننده آن آرگومان را ننویسد، مقدارِ پیش‌فرض به‌طورِ خودکار جایش می‌نشیند. به
این ترتیب یک روال هم حالتِ پرکاربرد را پوشش می‌دهد و هم حالتِ استثنا را،
بی‌آنکه به نسخهٔ جداگانه‌ای از روال نیاز باشد:

\begin{salamfa}
روال درود(نام : رشته, پیشوند : رشته = "سلام", بار : صحیح = 1):
    تکرار بار:
        سرچاپ پیشوند, نام
    پایان
پایان

روال ریشه:
    درود("سارا")              // پیشوند = "سلام"، بار = ۱
    درود("علی", "های")        // بار = ۱
    درود("مکس", "هی", 3)      // همه داده شده
پایان
\end{salamfa}

\begin{progout}
سلام سارا
های علی
هی مکس
هی مکس
هی مکس
\end{progout}

عبارتِ پیش‌فرض \emph{هنگامِ فراخوانی} و هر بار که آرگومان نوشته نشود ارزیابی
می‌شود؛ پس می‌تواند به پایدارها و دیگر نام‌های سطح‌بالا اشاره کند، اما نه به
پارامترهای دیگرِ همان روال:

\begin{salamfa}
پایا گام : صحیح = 5
روال افزایش(عدد : صحیح, مقدار : صحیح = گام) : صحیح:
    برگشت عدد + مقدار
پایان

روال ریشه:
    سرچاپ افزایش(10)       // مقدار = گام = ۵
    سرچاپ افزایش(10, 1)    // مقدار = ۱
پایان
\end{salamfa}

\begin{progout}
15
11
\end{progout}

\begin{warn}
پیش‌فرض‌ها باید در \emph{انتها} باشند: پس از نخستین پارامترِ دارای پیش‌فرض، همهٔ
پارامترهای بعدی هم باید پیش‌فرض داشته باشند. پس
\code{(الف صحیح, ب صحیح = 0)} درست است، اما
\code{(الف صحیح = 0, ب صحیح)} پذیرفته نمی‌شود اگر تنها یک آرگومان بدهید،
کامپایلر نمی‌تواند تشخیص دهد کدام پارامتر را در نظر داشته‌اید. مقدارِ پیش‌فرض هم
باید با گونهِ پارامتر سازگار باشد، پس \code{مقدار صحیح = "پنج"} خطای زمانِ
کامپایل است.
\end{warn}

\section{بازگشت خروجیِ روال}

واژهٔ \code{برگشت} (معادلِ \code{ret}) دو کار هم‌زمان انجام می‌دهد: مقداری را
به بیرونِ روال می‌فرستد، و بی‌درنگ از روال خارج می‌شود. در روالِ \code{بیشینه}،
اگر \code{الف > ب} باشد، \code{برگشت الف} اجرا می‌شود و باقیِ روال هرگز
اجرا نمی‌شود.

\begin{tip}
از همین ویژگیِ «خروجِ بی‌درنگ» می‌توان برای ساده‌تر کردنِ کد استفاده کرد.
روالِ \code{بیشینه} را ببینید: به‌جای نوشتنِ \code{اگر...وگرنه}، نخست حالتِ
خاص را با \code{برگشت} مدیریت می‌کنیم و باقی را در ادامه می‌نویسیم. این الگو
به «بازگشتِ زودهنگام» معروف است و کد را تخت‌تر و خواناتر می‌کند.
\end{tip}

\subsection{روال‌های بدونِ خروجی}

هر روالی لازم نیست چیزی برگرداند. روالی که فقط کاری انجام می‌دهد (مثلاً چیزی
چاپ می‌کند) گونهِ خروجی ندارد:

\begin{salamfa}
روال سلام بگو(نام : رشته):
    سرچاپ "سلام،", نام
    سرچاپ "به سلام خوش آمدی!"
پایان

روال ریشه:
    سلام بگو("سارا")
    سلام بگو("علی")
پایان
\end{salamfa}

دقت کنید که پس از \code{(نام رشته)} هیچ گونهِ خروجی نیامده چون این روال چیزی
برنمی‌گرداند. در این‌گونه روال‌ها، \code{برگشت} (بدونِ مقدار) را می‌توان برای
خروجِ زودهنگام به‌کار برد، اما اجباری نیست.

\section{بازگشتی: روالی که خودش را صدا می‌زند}

یکی از زیباترین ایده‌های برنامه‌نویسی این است که یک روال می‌تواند \emph{خودش}
را صدا بزند. به این کار \textbf{برگشت} یا \textbf{بازگشتی} (recursion)
می‌گویند. مثالِ کلاسیک، فاکتوریل است:

\begin{salamfa}
روال فاکتوریل(ن : صحیح) : صحیح:
    اگر ن <= 1:
        برگشت 1
    پایان
    برگشت ن * فاکتوریل(ن - 1)
پایان

روال ریشه:
    سرچاپ فاکتوریل(5)
    سرچاپ فاکتوریل(7)
پایان
\end{salamfa}

فاکتوریلِ ۵ یعنی $5\times4\times3\times2\times1 = 120$. روال چنین می‌اندیشد:
«فاکتوریلِ ۵ برابر است با ۵ ضربدرِ فاکتوریلِ ۴، و فاکتوریلِ ۴ برابر است با ۴
ضربدرِ فاکتوریلِ ۳، و\ldots» تا به فاکتوریلِ ۱ برسد که پاسخش روشن است: ۱.

\begin{warn}
هر روالِ بازگشتی به یک \textbf{شرطِ توقف} (پایهٔ بازگشت) نیاز دارد جایی که
دیگر خودش را صدا نمی‌زند. در فاکتوریل، این شرط \code{ن <= 1} است. اگر این شرط
را فراموش کنید، روال تا بی‌نهایت خودش را صدا می‌زند و برنامه از کار می‌افتد. این
خطا مانندِ همان «حلقهٔ بی‌پایانِ» فصلِ پیش است.
\end{warn}

\begin{deep}[نگاهی به درون: بازگشت چگونه کار می‌کند]
هر بار که روالی صدا زده می‌شود چه خودش، چه دیگری رایانه یک «قاب» تازه
روی \emph{پشتهٔ فراخوانی} (call stack) می‌سازد که پارامترها و متغیرهای آن
فراخوانی را نگه می‌دارد. در فاکتوریلِ ۵، پنج قاب روی هم چیده می‌شوند تا به پایهٔ
بازگشت برسیم؛ سپس یکی‌یکی باز می‌گردند و ضرب‌ها انجام می‌شود. به همین دلیل
بازگشتِ بیش‌از‌اندازه عمیق می‌تواند پشته را پُر کند نکته‌ای که در فصلِ ۲۰
به آن بازمی‌گردیم.
\end{deep}

\section{ترتیبِ تعریفِ روال‌ها}

در نمونه‌های این کتاب، روال‌های کمکی را معمولاً \emph{پیش} از روالِ \code{ریشه}
می‌نویسیم. اما این تنها یک قراردادِ زیباست، نه یک اجبار. سلام برنامه را در دو
گذر بررسی می‌کند: نخست همهٔ روال‌ها را \emph{می‌شناسد}، سپس بدنه‌ها را وارسی
می‌کند. به همین دلیل می‌توانید روالی را \emph{پیش} از تعریفش هم صدا بزنید ---
چیزی که گاهی «اعلانِ پیش‌رو» نامیده می‌شود و در زبان‌های دیگر دردسرساز است، اما
در سلام به‌طورِ طبیعی کار می‌کند.

\section{ترکیبِ روال‌ها}

زیباییِ روال‌ها آن‌جاست که خروجیِ یکی را می‌توان \emph{ورودیِ} دیگری کرد. در
مثالِ نخستِ این فصل دیدیم:

\begin{salamfa}
سرچاپ "بیشینه:", بیشینه(توان(2, 4), توان(3, 2))
\end{salamfa}

اینجا نخست \code{توان(2, 4)} برابرِ ۱۶ و \code{توان(3, 2)} برابرِ ۹ محاسبه
می‌شوند، سپس \code{بیشینه(16, 9)} برابرِ ۱۶ می‌شود. این «لانه‌سازیِ» روال‌ها، پایهٔ
ساختنِ برنامه‌های بزرگ از قطعاتِ کوچک است.

\section{خلاصهٔ فصل}

\begin{itemize}
  \item روال، بسته‌ای نام‌دار از دستورهاست که از تکرار جلوگیری می‌کند و کد را
        خوانا می‌سازد.
  \item پارامترها ورودیِ روال‌اند و هنگامِ صدا زدن مقدار می‌گیرند.
  \item \code{برگشت} خروجی را می‌فرستد و از روال خارج می‌شود.
  \item پارامترهای پایانی می‌توانند \emph{مقدارِ پیش‌فرض} داشته باشند
        (\code{ب صحیح = 0}) که فراخواننده می‌تواند آن را ننویسد.
  \item روال‌ها می‌توانند بدونِ خروجی باشند.
  \item روالِ بازگشتی خودش را صدا می‌زند و به شرطِ توقف نیاز دارد.
  \item خروجیِ یک روال می‌تواند ورودیِ روالِ دیگر باشد.
\end{itemize}

\section{تمرین‌ها}

\exercise روالی به نامِ \code{مربع} بنویسید که یک عدد بگیرد و مربعِ آن را
برگرداند.

\exercise روالی به نامِ \code{کمینه} بنویسید که دو عدد بگیرد و کوچک‌ترشان را
برگرداند.

\begin{sloppypar}
\exercise روالِ بازگشتیِ \code{فیبوناچی} را بنویسید
(\code{فیبوناچی(ن) = فیبوناچی(ن-۱) + فیبوناچی(ن-۲)}) با شرطِ توقفِ مناسب.
\end{sloppypar}

\exercise روالِ \code{توان(پایه صحیح, نما صحیح = 2) صحیح} را بنویسید که
\code{پایه} را به توانِ \code{نما} برساند. آن را به‌صورتِ \code{توان(5)} (باید ۲۵
شود) و \code{توان(2, 10)} (باید ۱۰۲۴ شود) صدا بزنید. سپس پارامتری \emph{بدونِ}
پیش‌فرض پس از \code{نما} اضافه کنید و خطای کامپایلر را بخوانید.

\exercise روالی بنویسید که یک عدد بگیرد و بگوید اول است یا نه (راهنمایی: با یک
حلقه بررسی کنید که آیا مقسوم‌علیهی غیر از ۱ و خودش دارد).

\exercise روالی به نامِ \code{ب.م.م} بنویسید که بزرگ‌ترین مقسوم‌علیهِ مشترکِ دو
عدد را با روشِ اقلیدس حساب کند.
